\section{Програмна реализация}

След представянето на изискванията, архитектурата и домейн модела следва анализ на конкретната програмна реализация. Тази глава разглежда последователно три аспекта на изграденото приложение: реализацията на back-end логиката чрез контролери, услуги и repository компоненти; представящия слой и потребителското взаимодействие чрез Thymeleaf шаблони; и операционната среда -- структурата на проекта, конфигурацията и build процеса. Целта е архитектурните идеи от предходната глава да бъдат проследени в тяхната програмна материализация.

\subsection{Реализация на back-end логиката}

След описването на архитектурата и домейн модела е необходимо да се разгледа конкретната имплементация на back-end логиката. В този раздел фокусът е върху действителното поведение на системата, реализирано чрез контролери, услуги, repository компоненти и конфигурация. Анализът следва структурата на проекта и показва как отделните модули взаимодействат, за да реализират регистрация, управление на фирми и продукти, работа с количка, checkout и история на поръчките.

За разлика от предходните глави, които се концентрират върху общите принципи, настоящият раздел стъпва директно върху наличния сорс код. Именно това прави изложението особено ценно: тук архитектурните идеи се виждат в тяхната програмна материализация.

\subsubsection{Конфигурация на основните инфраструктурни компоненти}

Back-end поведението на приложението зависи от няколко инфраструктурни компонента, дефинирани централизирано. Най-важните сред тях са \texttt{ModelMapper} и \texttt{PasswordEncoder}. Те се регистрират като Spring bean обекти и впоследствие се инжектират в service компонентите.

\begin{lstlisting}[caption={Конфигурация на mapper и password encoder},label={lst:bean-config}]
@Configuration
public class ApplicationBeanConfiguration {

    @Bean
    public ModelMapper modelMapper() {
        return new ModelMapper();
    }

    @Bean
    public PasswordEncoder passwordEncoder() {
        return new Pbkdf2PasswordEncoder();
    }
}
\end{lstlisting}

Листинг \ref{lst:bean-config} показва добър пример за централизирано управление на зависимости. \texttt{ModelMapper} подпомага преобразуването между различни модели, а \texttt{Pbkdf2PasswordEncoder} осигурява по-сигурно съхранение на паролите. Макар тази конфигурация да е кратка, тя има ключово значение за правилното функциониране на регистрацията и профилния модул.

\subsubsection{Потребителски модул}

Потребителският модул е най-обемен от гледна точка на уеб потоци. Той обхваща регистрация, потвърждение на акаунт, вход, изход, извличане на потребителска информация, редакция на профил и изчисляване на стойността на продуктите в количката.

\paragraph{Регистрация на потребител}

Регистрационният процес започва в \texttt{UserController}, където POST маршрутът за регистрация приема \texttt{UserRegisterBindingModel}. Ако входните данни са коректни от гледна точка на binding и няма вече съществуващ потребител със същия имейл, контролерът делегира работата към \texttt{UserServiceImpl}. Там се създава нов entity обект от тип \texttt{User}, инициализират се колекциите за компании, количка и история, задават се основните атрибути и паролата се хешира преди запис.

След запис на потребителя системата изпраща електронно писмо с линк за потвърждение. Това е съществена част от логиката, защото не позволява автоматичен вход преди потвърждаване на акаунта. От чисто инженерна гледна точка модулът демонстрира добра последователност: контролерът управлява потока, услугата реализира бизнес логиката, а repository слоят извършва persistence операцията.

\begin{lstlisting}[caption={Опростен вид на регистрационната логика},label={lst:user-register}]
@Override
public Optional<UserRegisterBindingModel> register(
        UserRegisterBindingModel input) {
    if (userRepository.findByEmail(input.getEmail()).isPresent()) {
        return Optional.empty();
    }

    User user = new User();
    user.setCompanies(new ArrayList<>());
    user.setProductsCart(new ArrayList<>());
    user.setHistoryProducts(new ArrayList<>());
    user.setFirstName(input.getFirstName());
    user.setLastName(input.getLastName());
    user.setEmail(input.getEmail());
    user.setConfirmed(false);
    user.setPassword(passwordEncoder.encode(input.getPassword()));
    user = userRepository.saveAndFlush(user);

    SimpleMailMessage message = new SimpleMailMessage();
    message.setTo(input.getEmail());
    message.setSubject("Email Confirmation");
    message.setText(
            "http://localhost:8080/users/"
                    + user.getId()
                    + "/confirm-email");
    emailSender.send(message);

    return Optional.of(input);
}
\end{lstlisting}

Листинг \ref{lst:user-register} илюстрира същината на регистрационния процес. В него ясно се вижда връзката между persistence слоя, сигурността на паролите и e-mail интеграцията. Същевременно могат да се забележат и някои ограничения: URL адресът за потвърждение е твърдо зададен към \texttt{localhost:8080}, а подателят и конфигурацията на пощата зависят от локални настройки.

\paragraph{Вход в системата}

Входът също е реализиран в \texttt{UserController}. Контролерът подава входния модел към потребителската услуга, която намира потребителя по имейл, проверява дали акаунтът е потвърден и сравнява въведената парола с хеша в базата. При успех имейлът се записва в \texttt{HttpSession}. Този подход е относително прост, но напълно достатъчен за демонстрационен проект.

Важно е да се отбележи, че логиката не използва Spring Security filter chain или security context. Вместо това достъпът до защитени ресурси се контролира ръчно чрез проверки за наличен атрибут в сесията. От методологична гледна точка това прави кода лесен за следване, но въвежда по-голям риск от дублиране на проверки и от пропуски при добавяне на нови маршрути.

\paragraph{Потвърждение на акаунт и управление на профил}

Потвърждението на акаунта е реализирано чрез маршрут \texttt{/users/\{userId\}/confirm-email}. При извикване на този маршрут потребителската услуга намира съответния потребител и маркира профила като потвърден. След това приложението връща екран за вход.

Профилният модул позволява извличане и обновяване на потребителските данни. Логиката включва две важни проверки: дали текущата сесия е активна и дали старата парола съвпада със съхранения хеш. По този начин системата реализира базов контрол върху промяната на чувствителна информация.

\subsubsection{Модул за фирми}

Управлението на фирмите е реализирано чрез \texttt{CompanyController} и \texttt{CompanyServiceImpl}. От функционална гледна точка този модул играе ролята на подготвителен етап за продуктовия каталог, защото всяка фирма може да бъде използвана при въвеждане на продукти.

Контролерът за фирми е относително компактен. Той осигурява GET маршрут за визуализация на формата и POST маршрут за запис на фирма. При POST обработката се използва \texttt{AddCompanyBindingModel}, който валидира полета като име, адрес, ДДС номер и регистрационен статус. Service имплементацията намира текущия потребител по имейл от сесията, създава entity обект \texttt{Company}, интерпретира текстовото поле за регистрационен признак и инициализира асоциираната колекция от продукти.

Тук се вижда интересна особеност на дизайна: фирмата не е абстрактна част от глобален каталог, а принадлежи към конкретен потребител. Това подсилва идеята, че приложението е по-близо до персонален търговски workspace, отколкото до многостранна публична маркетплейс система.

\subsubsection{Модул за продукти и каталог}

Продуктовият модул е ключова част от приложението, защото именно чрез него системата придобива своята каталожна функционалност. Този модул е разделен между \texttt{ProductController}, \texttt{ProductsService} и \texttt{ProductServiceImpl}. Освен това разчита на \texttt{CompanyRepository}, \texttt{ProductRepository}, \texttt{UserRepository} и \texttt{ShoppingCartService}.

\paragraph{Добавяне на продукт}

Добавянето на продукт започва с форма, която очаква име, цена, количество, категория, фирма и URL към изображение. Контролерът валидира входния модел и при успех извиква услугата за продукти. Там първо се намира фирмата по име, след което се създава нов обект \texttt{Product}. Категорията се преобразува чрез \texttt{CategoryEnum.valueOf()}, а останалите полета се попълват директно от входния модел.

\begin{lstlisting}[caption={Създаване на нов продукт в service слоя},label={lst:add-product-service}]
@Override
public void addProduct(
        ProductAddBindingModel input,
        Object email) throws Throwable {
    Company company = companyRepository.findByName(input.getCompany())
        .orElseThrow(() -> new Throwable("Company not found"));

    Product product = new Product();
    product.setUrl(input.getUrl());
    product.setCategory(CategoryEnum.valueOf(input.getCategory()));
    product.setCompany(company);
    product.setPrice(input.getPrice());
    product.setName(input.getName());
    product.setQuantity(input.getQuantity());

    productRepository.save(product);
}
\end{lstlisting}

Листинг \ref{lst:add-product-service} показва добре една типична service операция: извличане на зависим домейн обект, създаване на нов entity обект, попълване на полета и persistence чрез repository интерфейс. Решението е компактно и четимо, което е предимство за учебна разработка.

\paragraph{Каталог и филтриране по категории}

Каталогът е реализиран чрез маршрут от вида \texttt{/products/catalog/\{category\}}. В контролера се използва \texttt{switch} конструкция, която делегира към различни service методи според категорията. Тези service методи извличат всички продукти и филтрират резултатите по стойността на \texttt{CategoryEnum}. От архитектурна гледна точка това решение е достатъчно просто и разбираемо. То не е най-ефективният възможен подход, тъй като филтрирането се извършва на ниво Java поток върху колекцията от всички продукти, но за ограничен набор от данни е напълно приемливо.

Каталогът включва и клиентско филтриране по ценови диапазон, което използва REST крайната точка за извличане на максимална цена. Това е интересен пример за комбиниране на сървърно рендиране и допълнителна динамика от JavaScript.

\paragraph{Преглед на детайли и добавяне в количка}

За всеки продукт системата предоставя детайлен екран, от който потребителят може да заяви определено количество за добавяне в количката. Контролерът извиква метод за проверка на наличността, а при положителен резултат делегира към service логиката за реално добавяне на продукта.

В \texttt{ProductServiceImpl} това става чрез намаляване на наличността на продукта и добавяне или обновяване на запис в количката. Ако продуктът вече е намерен в количката, съществуващият запис се увеличава с новото количество; ако не съществува, създава се нов \texttt{ShoppingCartEntity}. Този модел е функционален, но тук се откроява и важна архитектурна слабост: проверката дали продуктът вече е в количката се извършва глобално върху всички записи в количката, а не само върху количката на текущия потребител. Освен това сравнението се базира на името на продукта, а не на комбинация от потребител и идентификатор на продукт. Това е достатъчно сериозно ограничение и следва да бъде изрично отбелязано.

\subsubsection{Количка и финализиране на поръчка}

Модулът за количка и checkout е реализиран чрез комбинация между \texttt{ProductController}, \texttt{ShoppingCartServiceImpl}, \texttt{UserServiceImpl} и \texttt{HistoryServiceImpl}. Това е една от най-интересните части на back-end логиката, защото координира няколко различни типа операции.

\paragraph{Преглед на количката}

При отваряне на количката контролерът добавя към модела списък от записи в количката, брой продукти, обща цена и крайна цена с доставка. Тук обаче отново е налице архитектурна особеност: методът \texttt{findAllShoppingCartEntities()} връща всички записи в количката, без филтриране по потребител. Това означава, че при реална многопотребителска среда визуализацията на количката не е строго изолирана на ниво потребител, което е съществено ограничение на текущата имплементация.

Въпреки това самият поток е ясно структуриран и лесен за проследяване. Потребителят вижда събраните продукти, цените и формата за адрес, а след това може да финализира поръчката чрез отделен POST маршрут.

\paragraph{Checkout логика и запис в историята}

Checkout процесът е реализиран в \texttt{HistoryServiceImpl} и е маркиран като \texttt{@Transactional}. Това е разумно решение, тъй като операцията променя множество състояния: създава нов исторически запис, обновява или изчиства количката на потребителя, изтрива определени записи и премахва продукти с нулева наличност.

\begin{lstlisting}[caption={Същност на checkout операцията},label={lst:checkout-service}]
@Override
@Transactional
public void secureCheckout(
        String address,
        Object email,
        Double totalPrice) throws Throwable {
    User user = userRepository.findByEmail(email.toString())
        .orElseThrow(() -> new Throwable("User not found"));

    HistoryEntity historyEntity = new HistoryEntity();
    historyEntity.setAddress(address);
    historyEntity.setPrice(totalPrice);
    historyEntity.setDate(LocalDate.now());
    historyEntity.setUser(user);
    historyRepository.save(historyEntity);

    setNullToShoppingCartProducts(user.getId());
    user.setProductsCart(new ArrayList<>());
    user.getHistoryProducts().add(historyEntity);
    userRepository.save(user);

    shoppingCartRepository.deleteAllByUser_Id(user.getId());
    deleteAllProductsWithZeroQuantity();
}
\end{lstlisting}

Листинг \ref{lst:checkout-service} е особено показателен. От една страна, той демонстрира добре как service слой може да координира сложна домейн операция. От друга страна, разкрива и опростения характер на модела. Вместо реална поръчка с отделни позиции се създава обобщен \texttt{HistoryEntity}. Освен това при приключване на покупка продуктите с нулева наличност се изтриват напълно, вместо просто да останат в каталога като изчерпани. Това е валиден компромис за учебен проект, но в реална система би било по-удачно продуктите да останат видими с индикация за липса на наличност.

\paragraph{Изчистване на историята}

Историята може да бъде изчиствана от потребителя. Реализацията последователно премахва асоциациите от страна на потребителя, занулява връзката към потребителя в историческите записи и изтрива вече осиротелите записи. Тази логика демонстрира осъзнаване на релационните зависимости, макар и реализирано чрез опростен процедурен подход.

\subsubsection{REST слой}

Освен HTML ориентираното поведение системата включва и \texttt{MaxPriceRestController}, който обслужва маршрут \texttt{/api/products/max-price}. Контролерът извиква \texttt{ProductsService.getProductMaxPrice()} и връща \texttt{MaxPriceViewModel} чрез \texttt{ResponseEntity}. Реализацията е минималистична, но добре подбрана. Тя показва, че същият домейн модел може да бъде използван и за машинно четим интерфейс.

От гледна точка на учебната стойност това е важно, защото демонстрира как в едно и също приложение могат да съществуват както класически MVC маршрути, така и REST крайни точки. В последствие именно тази крайна точка се използва от JavaScript логиката в каталога за конфигуриране на клиентския ценови филтър.

\subsubsection{Роля на repository интерфейсите в реализацията}

Макар repository интерфейсите да са относително кратки, те играят важна роля в back-end имплементацията. Чрез тях услугите извличат потребители по имейл, фирми по име, продукти по идентификатор, записи от количка по потребител и история на поръчки по потребителски идентификатор. Този стил на работа подчертава предимството на Spring Data JPA: много от ежедневните persistence операции могат да бъдат описани с кратки и четими методи.

Същевременно именно на това ниво се вижда и една от възможностите за бъдеща оптимизация. Част от филтрирането в момента се извършва в service слоя след извличане на големи колекции. При по-натоварена система би било по-ефективно част от тази логика да бъде преместена към по-целенасочени repository заявки.

\subsubsection{Наблюдавани силни страни на back-end реализацията}

На база направения анализ могат да бъдат формулирани няколко положителни характеристики на back-end слоя:

\begin{itemize}[leftmargin=*]
    \item последователно използване на service слой за по-съществените бизнес операции;
    \item отделяне на HTML контролери и REST контролер;
    \item използване на mapper и binding модели за по-добра структурираност на данните;
    \item прилагане на password hashing и e-mail confirmation механизъм;
    \item използване на транзакционност при операция, която променя множество зависими състояния.
\end{itemize}

Тези качества правят системата убедителна като учебен пример и осигуряват добра основа за бъдещо надграждане.

\subsubsection{Ограничения и технически рискове в back-end слоя}

Наред със силните страни back-end кодът съдържа и няколко важни ограничения, които трябва да бъдат отчетени честно:

\begin{itemize}[leftmargin=*]
    \item контролът на достъпа е реализиран ръчно чрез сесийни проверки в контролерите, без централизирана security архитектура;
    \item checkout, количка и проверката за вече добавен продукт работят с логика, която не е строго изолирана по потребител;
    \item в част от операциите се разчита на твърдо зададени стойности, като URL за потвърждение и конфигурация за локално изпълнение;
    \item продуктите с нулева наличност се изтриват, вместо да преминават в състояние „изчерпан“;
    \item част от домейн логиката използва по-опростени структури, подходящи за учебен контекст, но недостатъчни за зряла търговска система.
\end{itemize}

Тези ограничения не следва да се разглеждат като провал на проекта. Напротив, те показват къде е естествената граница между функционален демонстрационен прототип и production-ready система.

\subsection{Представящ слой и потребителско взаимодействие}

Back-end логиката на една система за електронна търговия има стойност само ако бъде експонирана чрез последователен и използваем потребителски интерфейс. В разглеждания проект визуалният слой е реализиран чрез Thymeleaf шаблони, допълнени от CSS, Bootstrap базирани компоненти и ограничено количество JavaScript. Този избор е в пълно съответствие с общата архитектурна постановка на приложението: интерфейсът е сървърно рендиран, което улеснява интеграцията между модели, валидация и HTML форми.

Настоящият раздел разглежда организацията на шаблоните, ролята на навигацията, локализацията, формите и основните UX решения. Анализът се базира върху реалните HTML шаблони в директорията \texttt{src/main/resources/templates} и статичните ресурси в \texttt{src/main/resources/static}.

\subsubsection{Организация на шаблоните}

Визуалният слой е структуриран около набор от отделни Thymeleaf шаблони, всеки от които отговаря на конкретен потребителски екран. Сред тях се открояват следните страници:

\begin{itemize}[leftmargin=*]
    \item \texttt{index.html} -- начална страница и входна точка за неавтентикирани потребители;
    \item \texttt{login.html} и \texttt{register.html} -- форми за удостоверяване и регистрация;
    \item \texttt{profile.html} -- екран за преглед и редакция на профила;
    \item \texttt{add-company.html} -- форма за добавяне на фирма;
    \item \texttt{add-product.html} -- форма за добавяне на продукт;
    \item \texttt{products-catalog.html} -- основен каталожен екран;
    \item \texttt{details.html} -- детайлен изглед на конкретен продукт;
    \item \texttt{shopping-cart.html} -- количка и форма за checkout;
    \item \texttt{history.html} -- история на направените поръчки;
    \item \texttt{users.html}, \texttt{companies.html}, \texttt{delivery.html} и error шаблоните -- допълнителни визуални компоненти на системата.
\end{itemize}

Този набор от шаблони е достатъчен, за да обхване целия потребителски цикъл в приложението. Важно е да се отбележи, че визуалният слой не е структуриран чрез централен layout механизъм или фрагменти, а по-скоро чрез отделни самостоятелни HTML файлове. Това прави проекта лесен за локално проследяване, но води и до значително дублиране на елементи като навигационната лента.

\begin{table}[H]
    \centering
    \caption{Съответствие между шаблони и основни потребителски функции}
    \label{tab:template-mapping}
    \begin{tabularx}{\textwidth}{L{3.3cm}YY}
        \toprule
        Шаблон & Основна цел & Потребителски контекст \\
        \midrule
        \texttt{index.html} & Представяне на начална страница и достъп до основни действия & Гост или неавтентикиран потребител \\
        \texttt{register.html} & Създаване на нов акаунт & Гост \\
        \texttt{login.html} & Вход в системата & Гост или непотвърден потребител \\
        \texttt{profile.html} & Преглед и редакция на данни, списък с фирми & Удостоверен потребител \\
        \texttt{add-company.html} & Въвеждане на фирма & Удостоверен потребител \\
        \texttt{add-product.html} & Въвеждане на продукт & Удостоверен потребител \\
        \texttt{products-catalog.html} & Каталожен преглед и филтриране & Удостоверен потребител \\
        \texttt{details.html} & Избор на конкретен продукт и количество & Удостоверен потребител \\
        \texttt{shopping-cart.html} & Обобщение на количка и checkout & Удостоверен потребител \\
        \texttt{history.html} & Преглед и изчистване на историята & Удостоверен потребител \\
        \bottomrule
    \end{tabularx}
\end{table}

Таблица \ref{tab:template-mapping} показва, че интерфейсът е построен около ясни и тематично отделени екрани. Това е добра практика, защото всяка страница има разпознаваема функция и лесно може да бъде обвързана с конкретен контролерен маршрут.

\subsubsection{Навигация и общи интерфейсни елементи}

Един от най-видимите повторяеми елементи в шаблоните е навигационната лента. Тя присъства в почти всички екрани и включва връзки към регистрация, вход, профил, каталог, добавяне на продукт, история, количка, изход и смяна на езика. От гледна точка на потребителската логика това е силно решение: основните действия са достъпни от почти всеки екран и навигационната структура остава позната.

В същото време архитектурно решението има недостатък. Понеже навигацията се повтаря ръчно в множество шаблони, всяка бъдеща промяна трябва да бъде приложена на много места. При по-голямо приложение това би увеличило вероятността от несъответствия и би направило поддръжката по-трудна. По-зрял подход би бил използването на Thymeleaf fragments или общ layout шаблон. Независимо от това, за настоящия проект текущият вариант е лесен за проследяване и не затруднява разбирането на страниците.

Допълнителен общ елемент е визуалният фон, който в много страници използва фиксирано изображение и сходна цветова схема. Това придава визуална последователност, макар и реализирана чрез inline стилове, а не чрез централен CSS дизайн слой.

\subsubsection{Thymeleaf като механизъм за свързване на данни и интерфейс}

Thymeleaf \cite{thymeleaf} играе централна роля в представящия слой. Чрез атрибути като \texttt{th:text}, \texttt{th:field}, \texttt{th:each}, \texttt{th:if}, \texttt{th:href} и \texttt{th:action} шаблоните могат директно да използват данните, подготвени от контролерите. Това значително улеснява връзката между back-end слоя и HTML формите.

Особено полезно е използването на \texttt{th:field} при формите. По този начин се постига естествено свързване между binding модел и HTML полета, както и възможност за визуализиране на състояние при грешка. Например при формите за добавяне на фирма, продукт или редакция на профил, полетата се обвързват директно с модела, а условната визуализация на грешки използва \texttt{\#fields.hasErrors(...)}.

\begin{lstlisting}[caption={Типичен Thymeleaf фрагмент за форма},label={lst:thymeleaf-form}]
<form th:action="@{/products/add-product}"
      th:object="${productAddBindingModel}"
      th:method="POST">
    <label for="name">Name</label>
    <input th:field="*{name}" id="name" type="text" class="form-control">
    <p th:if="${#fields.hasErrors('name')}" class="errors alert alert-danger">
        Name must be between 5 and 20 characters.
    </p>
</form>
\end{lstlisting}

Листинг \ref{lst:thymeleaf-form} показва как се реализира типичен формов сценарий. От една страна, HTML остава достатъчно четим. От друга страна, връзката с Java модела е непосредствена. Именно в това се състои едно от основните предимства на сървърно рендирания подход при Spring MVC приложения.

\subsubsection{Форми и обработка на потребителски вход}

Почти всички съществени действия в системата са реализирани чрез форми. Това включва регистрация, вход, редакция на профил, добавяне на фирма, добавяне на продукт, добавяне на количество в количката и checkout процеса. Формите са реализирани в стил, характерен за Bootstrap \cite{bootstrap}, с класове като \texttt{form-group}, \texttt{form-control}, \texttt{btn} и други визуални помощници.

В профилния екран се вижда сравнително богата композиция: форма за лични данни, визуализиране на текущите стойности и таблица с компаниите на потребителя. Това е добър пример за страница, която съчетава редакция и преглед на свързана информация. Формата за добавяне на фирма пък е по-компактна и логически разделя текстовите от числовите полета. Формата за добавяне на продукт е най-добре структурирана по отношение на домейнните данни, тъй като включва избор от enum категории и избор на фирма от предварително зареден списък.

От UX гледна точка позитивен аспект е, че при невалиден вход потребителят получава визуална обратна връзка чрез alert елементи. Това подобрява използваемостта и прави формовия поток по-ясен. Ограничение е, че съобщенията не са напълно стандартизирани на всички места, а част от тях са твърдо кодирани в самите HTML файлове.

\subsubsection{Каталог и продуктови детайли}

Каталогът е един от най-характерните екрани в системата. Той е реализиран като комбинация от навигационна зона, странично меню с категории, ценови филтър и галерия от продуктови карти. Визуалният модел се опира на Bootstrap, но включва и значително количество inline CSS, описващ оформлението на елементите, размерите на картите, поведението на менюто и отзивчивостта на страницата.

Съществено UX качество на каталога е, че всеки продукт се визуализира с изображение, име, категория и цена, а бутонът за детайлен преглед отвежда потребителя към следващата стъпка от взаимодействието. В допълнение клиентският филтър по цена придава усещане за по-динамичен интерфейс, без да се изисква пълноценна SPA архитектура.

Страницата с детайли на продукта е по-концентрирана върху единичен обект. Тя представя изображение, име, цена и наличност, а формата за въвеждане на количество позволява директно добавяне към количката. От гледна точка на потребителския поток това е правилно решение: каталогът служи за избор, а страницата за детайли -- за вземане на решение и заявяване на покупка.

\subsubsection{Количка и checkout екран}

Екранът за количка комбинира две взаимосвързани цели: визуализация на текущото съдържание и подготовка на финализирането на поръчката. Визуалната композиция разделя съдържанието на две части: списък с избраните продукти и summary панел със суми, доставка, адрес и бутон за checkout.

Това решение е потребителски удачно, тъй като показва едновременно детайла на избраните артикули и обобщената финансова рамка на поръчката. Формата за адрес е вградена директно в summary панела, което намалява броя на преходните екрани. От друга страна, самата страница е относително тежка като inline CSS и в production контекст би било по-удачно стиловете да се централизирани във външен ресурс.

\subsubsection{Профил и история на поръчките}

Профилният екран е сред най-богатите откъм комбинирано съдържание. Той събира както form-based редакция на основните потребителски полета, така и индиректно насърчава управлението на фирмите чрез страничен панел със списък от съществуващи компании и връзка за добавяне на нова. Това е добра UX практика, тъй като поставя логически свързаните операции в един и същ контекст.

Историята на поръчките е реализирана с тъмен табличен интерфейс, който се различава визуално от останалите екрани. Това прави страницата разпознаваема и внушава архивен характер на информацията. Таблицата показва идентификатор на поръчката, имейл, адрес, дата и цена. Допълнително е наличен бутон за изчистване на историята. От потребителска гледна точка този екран изпълнява коректно ролята си, макар че в production система биха били желани филтриране, търсене и по-богат контекст за отделните поръчки.

\subsubsection{Локализация и двуезичен интерфейс}

Поддръжката на два езика е реализирана чрез \texttt{messages.properties} и \texttt{messages\_bg.properties}, заедно с \texttt{LocaleChangeInterceptor}. В навигационната лента са добавени преки връзки \texttt{?lang=en} и \texttt{?lang=bg}, които позволяват смяна на езика чрез параметър в заявката. Това е прост, но ефективен модел за локализация в сървърно рендирана система.

От анализа на шаблоните се вижда, че голяма част от основните надписи използват message ключове чрез \texttt{\#\{...\}}. Това е добро решение, защото държи езиковите ресурси извън самите шаблони. Съществуват обаче и частично локализирани или напълно твърдо кодирани елементи. Например част от текстовете за категории, някои помощни изрази и някои съобщения за грешка не са напълно синхронизирани между двата езика. Това означава, че локализацията е реално налична, но все още не е напълно завършена.

\subsubsection{Използване на външни библиотеки и клиентски ресурси}

Визуалният слой използва Bootstrap, Font Awesome, jQuery и jQuery UI \cite{bootstrap,jquery}. Това е видно както от шаблоните за каталог и детайли, така и от количеството CDN препратки в HTML файловете. От гледна точка на бързото разработване това е разбираем избор: готовите UI библиотеки позволяват по-бързо постигане на визуално приемлив резултат.

Наред с положителните страни тук има и няколко наблюдения. Първо, в различни шаблони се използват различни версии на Bootstrap. Второ, част от ресурсите са включени едновременно от локални файлове и от CDN. Трето, голямо количество стилове се дефинира inline във всеки отделен шаблон, вместо да се поддържа по-строга централизирана front-end организация. Това не е критичен проблем за малък проект, но е важен индикатор за бъдещи възможности за подобрение.

\subsubsection{UX силни страни}

Въпреки опростения си характер, интерфейсът има няколко съществени предимства:

\begin{itemize}[leftmargin=*]
    \item основните действия са видими и достъпни от навигацията;
    \item формите са ясно отделени по функции и показват грешки при входни проблеми;
    \item каталогът е визуално разпознаваем и поддържа филтриране по категория и цена;
    \item checkout екранът обединява продуктова информация и обобщение на поръчката;
    \item профилът и историята предлагат логично групирани потребителски сценарии.
\end{itemize}

Тези характеристики правят системата използваема и подходяща за демонстрация пред научен ръководител, комисия или крайни потребители в контекста на дипломната разработка.

\subsubsection{Ограничения на представящия слой}

Както и при back-end частта, представящият слой има своите ограничения:

\begin{itemize}[leftmargin=*]
    \item дублиране на навигационна логика и HTML блокове между шаблоните;
    \item смесване на външни CDN ресурси, локални файлове и inline CSS;
    \item частична, но не пълна локализация на интерфейса;
    \item отсъствие на централен layout или fragment система;
    \item наличие на твърдо кодирани визуални и текстови елементи, които намаляват гъвкавостта.
\end{itemize}

Тези ограничения обаче не отменят факта, че интерфейсът е функционален и достатъчно изразителен за целите на проекта. По-скоро те очертават естествена траектория за последващо професионализиране на front-end слоя.

\subsection{Структура на проекта, конфигурация и експлоатационна среда}

Освен архитектурата и домейн логиката, всяко завършено приложение следва да бъде разглеждано и от operational перспектива: как е организирано в хранилището, какви зависимости използва, как се компилира, къде се съхраняват ресурсите и какви артефакти се генерират в резултат от build процеса. За настоящата система това е особено важно, тъй като проектът използва класическа Maven организация, комбинирана със Spring Boot конфигурация, MySQL свързаност, SMTP настройка и набор от HTML ресурси.

Този раздел анализира именно този аспект на разработката и показва как логическата структура на приложението се материализира във файловата организация на проекта.

\subsubsection{Обща организация на проекта}

Проектът е структуриран като стандартно Maven приложение. В кореновата директория се намират \texttt{pom.xml}, Spring Boot wrapper файловете, директорията \texttt{src} със сорс код, директорията \texttt{target} с генерирани артефакти, както и директорията \texttt{doc}, в която е разположена настоящата дипломна документация. Това разделение е важно, защото позволява да се разграничат поне три различни аспекта на системата:

\begin{itemize}[leftmargin=*]
    \item \textbf{изходен код и ресурси} -- реализирани в \texttt{src};
    \item \textbf{резултат от компилация и копиране на ресурси} -- реализирани в \texttt{target};
    \item \textbf{академична документация} -- реализирана в \texttt{doc}.
\end{itemize}

Тази подредба е полезна не само от практическа гледна точка, но и за дипломното описание, тъй като ясно отделя инженерния продукт от неговото документиране.

\begin{table}[H]
    \centering
    \caption{Ключови директории в проекта и тяхното предназначение}
    \label{tab:project-directories}
    \begin{tabularx}{\textwidth}{L{3.5cm}Y}
        \toprule
        Директория или файл & Предназначение \\
        \midrule
        \texttt{pom.xml} & Описва зависимостите, версията на Java и build конфигурацията \\
        \texttt{src/main/java} & Съдържа основния Java код на приложението \\
        \texttt{src/main/resources/templates} & Съдържа Thymeleaf шаблоните \\
        \texttt{src/main/resources/static} & Съдържа статични ресурси като CSS и JavaScript \\
        \texttt{src/main/resources/application.properties} & Съдържа локалната конфигурация за база данни, mail и JPA \\
        \texttt{target/classes} & Съдържа компилирани класове и копирани ресурси след build \\
        \texttt{target/test-classes} & Съдържа тестови артефакти при наличие на компилирани тестове \\
        \texttt{doc} & Съдържа LaTeX документацията и библиографията на проекта \\
        \bottomrule
    \end{tabularx}
\end{table}

Както се вижда от таблица \ref{tab:project-directories}, проектът следва добре позната и разбираема структура. Това е съществено предимство, защото разработчик, който е работил със Spring Boot и Maven, може бързо да се ориентира в него.

\subsubsection{Описание на зависимостите в \texttt{pom.xml}}

Файлът \texttt{pom.xml} е основният декларативен център на Java проекта. В него са зададени родителският Spring Boot артефакт, версията на Java и списъкът със зависимости. Анализът на файла показва, че системата използва сравнително типичен набор от библиотеки за Spring базирано уеб приложение.

\begin{table}[H]
    \centering
    \caption{Основни зависимости и роля в проекта}
    \label{tab:maven-dependencies}
    \begin{tabularx}{\textwidth}{L{4.5cm}YY}
        \toprule
        Зависимост & Функция в приложението & Практическо значение \\
        \midrule
        \texttt{spring-boot-starter-web} & Уеб инфраструктура и MVC маршрути & Осигурява основата на HTTP слоя \cite{springboot,springmvc} \\
        \texttt{spring-boot-starter-thymeleaf} & Сървърно рендирани HTML шаблони & Интегрира визуалния слой с Spring MVC \cite{thymeleaf} \\
        \texttt{spring-boot-starter-data-jpa} & Persistence слой чрез JPA & Улеснява работата с релационната база \cite{springdatajpa} \\
        \texttt{mysql-connector-java} & JDBC драйвер за MySQL & Позволява реална връзка с базата данни \cite{mysql} \\
        \texttt{spring-boot-starter-validation} & Валидация на входни данни & Използва се в binding моделите \cite{hibernatevalidator} \\
        \texttt{modelmapper} & Преобразуване между модели & Намалява ръчния код при mapping \cite{modelmapper} \\
        \texttt{spring-context-support} и \texttt{javax.mail} & Поддръжка за e-mail интеграция & Използват се при потвърждение на регистрацията \cite{springmail} \\
        \texttt{spring-security-crypto} & Хеширане на пароли & Осигурява \texttt{Pbkdf2PasswordEncoder} \cite{springsecuritycrypto} \\
        \texttt{spring-boot-starter-test} & Тестова инфраструктура & Осигурява база за бъдещи тестове \\
        \bottomrule
    \end{tabularx}
\end{table}

Таблица \ref{tab:maven-dependencies} показва, че изборът на зависимости е в общия случай добре съобразен с целите на проекта. Той покрива всички ключови нужди: уеб слой, шаблони, persistence, валидация, mail интеграция и password hashing.

Въпреки това анализът на \texttt{pom.xml} разкрива и няколко важни наблюдения. На първо място, зависимостта \texttt{spring-security-crypto} присъства два пъти с различни версии. На второ място, \texttt{spring-boot-starter-data-jpa} е декларирана с версия \texttt{LATEST}, което не е добра практика в устойчив проект, защото може да доведе до по-трудно възпроизводими build резултати. На трето място, някои security зависимости са коментирани, което подсказва, че проектът първоначално е бил разглеждан и като кандидат за по-богата security интеграция, но впоследствие е останал при по-опростен модел.

\subsubsection{Версия на платформата и езика}

Проектът използва Spring Boot 2.6.2 и Java 11. Това е разумна комбинация за учебно-приложна разработка. Java 11 е дългосрочно поддържана версия на езика и осигурява стабилна основа, а Spring Boot 2.6.2 предлага достатъчно зрял и добре документиран стек за уеб приложения. Тази комбинация гарантира съвместимост с използваните библиотеки и прави проекта лесен за стартиране в типична локална разработваческа среда.

\subsubsection{Конфигурация на приложението}

Операционното поведение на приложението се управлява основно от файла \texttt{src/main/resources/application.properties}. В него са зададени параметрите за логване, връзка към MySQL, JPA поведението, както и SMTP конфигурацията за електронната поща.

Конфигурацията на базата данни използва MySQL драйвер и URL към база \texttt{onlineShop}, като е включена опция за автоматично създаване на базата при нужда. Hibernate е настроен с \texttt{ddl-auto=update}, което означава, че схемата може да бъде синхронизирана с entity модела при стартиране. Това решение е удобно за локална разработка, защото намалява необходимостта от отделни миграционни скриптове в ранна фаза на проекта.

От operational гледна точка е важно и използването на \texttt{spring.jpa.open-in-view=false}. Това показва съзнателен избор да не се оставя persistence контекстът отворен през визуалния слой. Макар проектът да не използва сложни lazy loading сценарии в интерфейса, тази настройка е добра практика, защото намалява вероятността от неочаквани достъпи до базата по време на визуализация.

\subsubsection{Конфигурация на електронната поща}

Файлът \texttt{application.properties} съдържа и настройките за \texttt{spring.mail.host}, порт, SMTP автентикация, TLS и потребителско име. Тези параметри позволяват изпращане на потвърждаващ e-mail към новорегистрирани потребители. От функционална гледна точка това е значимо предимство, защото системата преминава отвъд чисто локалните CRUD операции и демонстрира интеграция с външен инфраструктурен компонент.

Наред с това обаче тук се вижда и едно от най-важните operational ограничения на проекта: чувствителни данни като парола за e-mail акаунт и данни за достъп до базата са записани директно в конфигурационен файл. Това е приемливо за локална учебна среда, но не следва да бъде допускано в production режим. По-зрял подход би включвал използване на environment variables, секретно хранилище или отделни профили за различни среди.

\subsubsection{Build процес и Maven жизнен цикъл}

Maven \cite{maven} осигурява стандартизиран жизнен цикъл за компилация, тестване и пакетиране. Дори когато проектът не използва сложни custom plugin сценарии, самото наличие на Maven build носи няколко преимущества:

\begin{itemize}[leftmargin=*]
    \item централизирано управление на зависимостите;
    \item възпроизводим build процес на различни машини;
    \item ясна конвенция за разположение на сорс код и ресурси;
    \item лесна интеграция със Spring Boot Maven plugin за стартиране и пакетиране.
\end{itemize}

В рамките на настоящия проект build процесът води до попълване на директорията \texttt{target}, където се намират \texttt{classes}, \texttt{test-classes}, както и стандартни поддиректории за генерирани сорсове. Най-важната от тях е \texttt{target/classes}, в която се копират както компилираните \texttt{.class} файлове, така и ресурсите от \texttt{src/main/resources}, включително HTML шаблоните.

\subsubsection{Съдържание и роля на \texttt{target}}

От наличната структура се вижда, че \texttt{target} съдържа поне следните поддиректории: \texttt{classes}, \texttt{generated-sources}, \texttt{generated-test-sources} и \texttt{test-classes}. Това е очакван резултат от Maven build процес. Практическото значение на тази директория е двойно.

Първо, тя съдържа реално изпълнимите артефакти на приложението след компилация. Второ, тя предоставя добра възможност за проверка дали ресурсите от \texttt{templates} и \texttt{static} са коректно копирани и достъпни за изпълнимата среда. В конкретния проект например шаблоните от \texttt{src/main/resources/templates} са налични и в \texttt{target/classes/templates}, което потвърждава правилното участие на ресурсите в build процеса.

\subsubsection{Практическа среда за изпълнение}

За локално стартиране на приложението са необходими поне следните предпоставки:

\begin{itemize}[leftmargin=*]
    \item налична Java 11 среда;
    \item работещ MySQL сървър;
    \item валидни данни за достъп до базата, съобразени с \texttt{application.properties};
    \item достъпна e-mail конфигурация, ако се тества регистрационният поток;
    \item Maven или използване на включения wrapper за стартиране на проекта.
\end{itemize}

Тази operational рамка е характерна за Spring Boot приложения и не изисква специализирана инфраструктура. Именно това е още едно предимство на избрания стек: проектът може да бъде разработван и демонстриран в сравнително стандартна среда без значителни допълнителни инструменти.

\subsubsection{Операционни рискове и възможности за подобрение}

Анализът на конфигурацията позволява да бъдат формулирани няколко operational риска:

\begin{itemize}[leftmargin=*]
    \item чувствителни параметри са съхранени в изходния код вместо във външна конфигурация;
    \item използването на \texttt{LATEST} в зависимост създава риск за повторяемостта на build процеса;
    \item наличието на дублирани зависимости подсказва необходимост от по-строга dependency хигиена;
    \item липсват ясно разграничени профили за development и production среда;
    \item отсъстват миграции на схемата чрез инструмент като Flyway или Liquibase.
\end{itemize}

Въпреки това operational картината на системата остава напълно достатъчна за целите на дипломната разработка. Проектът е стартиращ, конфигурируем и използва стандартни Java инструменти за компилация и изпълнение. Това го прави едновременно практичен и методологично удобен за анализ.

\subsubsection{Обобщение на програмната реализация}

Програмната реализация на проекта демонстрира ясно приложени принципи на Spring базирана уеб система: контролери за управление на HTTP потока, service слой за домейн логика, repository интерфейси за persistence, mapper за трансформации и инфраструктурна конфигурация за повторно използваеми компоненти. Представящият слой допълва тази организация чрез Thymeleaf шаблони и сървърно рендирани форми. Структурата на проекта, Maven конфигурацията и operational средата осигуряват добра основа за локално стартиране и анализ.

Същевременно направеният преглед показва и конкретни области за подобрение -- най-вече по отношение на сигурността, изолацията на потребителските данни, централизацията на front-end слоя и управлението на чувствителна конфигурация. Точно тази комбинация от работеща функционалност и ясно видими точки за развитие придава на проекта неговата инженерна стойност. Следващата глава представя ръководство за използване на системата чрез поредица потребителски сценарии и верификация на функционалността.
